\section{Visualisation Primitives}
\label{sec:foundations:primitives}

A visualisation primitive is a reusable chart form with a fixed visual
grammar: a defined set of axes, marks, and encodings that together
support a particular kind of analytical reading. The catalogue below
introduces the primitives drawn on later in the report. Each primitive
is described in terms of what it is, when it is used, and what kind of
analytical claim it supports; the specific choice of which primitive
serves which question in the proposed system is deferred to Chapter
\ref{chapter:solution}.

\subsection{Treemap}

A treemap is a space-filling diagram in which rectangular tiles encode
a hierarchy and tile area encodes a numeric measure. Tiles are nested:
the children of a category occupy the area of that category's tile,
recursively. The reading order is top-down for hierarchy and by area
for magnitude, so that the relative contribution of each subcategory to
its parent is immediately visible. Treemaps are cited by \cite{liu2022}
as one of the canonical density-by-category primitives.

\subsection{Calendar heatmap}

A calendar heatmap is a grid in which one axis is a calendar dimension
(day of week, hour of day, or date) and the other is a category or
another calendar dimension; cells are coloured by a numeric measure
according to a continuous colour scale. The primitive is used to expose
temporal patterns including weekly cycles, daily cycles, and seasonal
trends, because rhythmic structure in the underlying measure becomes
visible as banded patterns of similar colours.

\subsection{Parallel coordinates}

Parallel coordinates is a multi-dimensional visualisation in which each
vertical axis represents one dimension of the dataset and each
multi-dimensional record is rendered as a polyline that crosses every
axis at the height corresponding to its value on that axis. The
primitive is used to expose relationships among more than two
dimensions simultaneously and to support brushing across axes, so that
restricting the range on one axis isolates the polylines that fall
within that range across all the others. The primitive is cited by
\cite{liu2022} as one of the central tools of multi-dimensional
exploration.

\subsection{Hexagonal-bin grid}

A hexagonal-bin grid is a spatial-aggregation primitive that partitions
a planar area into hexagonal cells and colours each cell by an
aggregated measure, typically a count of events that fall within the
cell. Hexagonal cells are preferred over square bins because every
neighbouring cell is at the same centre-to-centre distance, which
produces a visually consistent appearance at varying zoom levels and
avoids the directional artefacts of square tessellation. The H3 system
of \cite{h3} provides a globally consistent hexagonal hierarchical
index whose cells can be referenced by stable identifiers across
resolutions.

\subsection{Scatter quadrant}

A scatter quadrant is a two-dimensional scatter plot with axes
deliberately chosen so that the four quadrants produced by drawing
reference lines at meaningful values correspond to qualitatively
distinct categories. For example, axes encoding cost and yield produce
four quadrants of low cost low yield, low cost high yield, high cost
low yield, and high cost high yield. Annotated quadrant boundaries
support categorical reading, and bubble size can be used to encode a
fourth measure beyond the two coordinate axes and any colour encoding.

\subsection{Sankey diagram}

A Sankey diagram is a flow diagram in which the width of each ribbon is
proportional to a flow magnitude. The diagram is read left to right (or
top to bottom): source categories on one side are connected by ribbons
to destination categories on the other, and the conserved width of the
ribbon across the diagram makes it easy to follow what fraction of a
particular source flows to a particular destination. The primitive is
used when the total flow between sources and destinations is itself a
quantity of analytical interest.

\subsection{Ranked bar}

A ranked bar chart is a horizontal bar chart whose bars are sorted by
the encoded numeric measure. It is the simplest primitive for surfacing
the top-$k$ of a distribution and for supporting direct numeric
comparison of category values, because the eye reads horizontal lengths
more accurately than it reads angles or areas. The primitive scales
gracefully to long category lists by limiting the rendered set to the
top-$k$ and providing a means of paging through the remainder.

\subsection{Two-dimensional projection scatter}

A two-dimensional projection scatter is a scatter plot whose axes are
derived by dimensionality reduction from a higher-dimensional feature
space. Each point represents a record described by many features; the
two coordinates are produced by an algorithm that attempts to preserve
some structural property of the high-dimensional space. The Uniform
Manifold Approximation and Projection (UMAP) algorithm of \cite{umap}
is a common choice that preserves local neighbourhood structure better
than alternatives such as t-distributed Stochastic Neighbour Embedding
(t-SNE) or principal component analysis for this purpose. The primitive
is used to expose clusters in a multi-feature dataset where no two
features alone separate the records cleanly.

\subsection{Time-distance trip diagram}

A time-distance trip diagram, also called a Marey diagram after its
nineteenth-century originator, places wall-clock time on the horizontal
axis and the ordinal position of each stop along a route on the
vertical axis. Each trip operated on the route is rendered as a
polyline whose slope reflects vehicle speed and whose vertical
discontinuities reflect dwell time at each stop. Bunching appears as
converging polylines, and gaps in service appear as wide vertical
voids in the diagram. The form originates with the train-scheduling
nomograms of Marey and has been extended to public transit by Lin and
Zeng \cite{linzeng2014}.

\subsection{Set-intersection matrix}

A set-intersection matrix is a matrix visualisation that exposes
intersections among set memberships across a fixed collection of sets.
Each column corresponds to one possible intersection of sets, marked by
filled dots in the rows of the sets it contains; a bar above each
column encodes the number of records that belong to that intersection,
and a bar beside each row encodes the total cardinality of that set.
The primitive is used when categorical membership is multi-valued, so
that a single record can belong to several sets at once. The form is
cited by \cite{liu2022} in its GroupSet variant.
